\chapter{Commissioning}
\label{chap:commissioning}

The task of this chapter is to show how the implemented solution was
brought into operation. For this project, this means the preparation of the
runtime environment, the execution of the experiment suite, and the steps
taken to make the results reproducible.

\section{Preparation of the Environment}
\label{sec:environment}

The experiments were carried out on a standard Linux machine with a CPU
only installation. The environment consists of Python~3.12 with the
packages PyTorch~2.8, torchvision, NumPy and Matplotlib. The MNIST raw
files were downloaded once from the official mirror and stored in the
project folder, so that the experiments can run without further downloads.
The installation and execution requires only the following commands:

\begin{lstlisting}[language=bash, caption={Installation and execution of
the experiment suite.}, label={lst:commands}]
pip install torch torchvision numpy matplotlib
cd moe_continual_learning/src
python3 run_experiments.py   # main experiment suite
python3 run_ablation.py      # load balancing ablation
python3 run_v2.py            # input level architecture variant
python3 make_figures.py      # figures for the report
python3 make_tables.py       # LaTeX tables for the report
\end{lstlisting}

\section{Execution of the Experiments}
\label{sec:execution}

The main experiment driver runs the full comparison matrix: four approaches
(naive MLP, MLP with EWC, wide MLP, MoE), two benchmarks (SplitMNIST and
PermutedMNIST) and three random seeds, which results in 24 complete
continual training runs. Before the construction of every model, the global
random seed is reset, so that all approaches start from comparable initial
conditions; the task permutations of PermutedMNIST are also derived from
the seed. After every training stage, the current accuracy row is printed
and stored, so the progress can be followed in the log file. A typical log
line looks as follows:

\begin{lstlisting}[caption={Example lines of the experiment log.},
label={lst:log}]
[split_mnist] seed=0 MoE: avg_acc=0.9523 forgetting=0.0561 bwt=-0.0561
[permuted_mnist] seed=2 MLP + EWC: avg_acc=0.9490 forgetting=0.0354
\end{lstlisting}

The results are written to disk after every single run. This proved to be
important in practice: during the project, one experiment process was
terminated by the operating system because of high memory usage, and thanks
to the incremental storage, no completed run had to be repeated. After this
incident, the PermutedMNIST loader was changed to apply the permutation on
the fly instead of materialising a copy of the whole data set per task,
which reduced the memory consumption of the benchmark clearly.

The runtime of a single continual training run over five tasks lies between
$18\,\mathrm{s}$ (naive MLP on SplitMNIST) and $430\,\mathrm{s}$ (MoE on
PermutedMNIST) on the CPU. The complete main suite takes about one hour,
the ablation about 20 minutes and the architecture variant about 25
minutes.

\section{Reproducibility Measures}
\label{sec:reproducibility}

Several measures were taken so that every number in this report can be
traced back to a concrete training run. All random choices are controlled
by seeds: the initialisation of every model, the data order of every
training run and the pixel permutations of PermutedMNIST. As a consequence,
repeated executions of the suite produce identical results, which was
verified explicitly when the suite had to be restarted after the incident
described above. The raw results of every run are stored as JSON files, and
the router profiles are stored separately as compressed NumPy archives.
Finally, all figures and all result tables of Chapter~7 were generated
directly from these stored results by the two scripts
\texttt{make\_figures.py} and \texttt{make\_tables.py}; no number was
copied by hand.
